\section{Sharp orders for the number of independent measurements}
\label{sec:sampling}

The preceding total-variation calculation concerns one observation. Repeated
observations require a separate argument: multiplying the one-observation
distance by the number of samples generally loses the correct scaling.
Here that gap can be closed for the stated, fully calibrated detector models.
All limits in this section hold as $s\downarrow0$, with detector parameters
fixed. Define the squared Hellinger distance and the affinity by
\[
 h^2(P,Q)=\int(\sqrt p-\sqrt q)^2,\qquad
 \rho(P,Q)=\int\sqrt{pq}=1-\frac{h^2(P,Q)}2.
\]
Integrals are replaced by sums for count distributions. This convention gives
$0\le h^2\le2$.

\subsection{Common leading moments at every fixed order}

The null condition and the uniform bound in \eqref{eq:Bp} imply
\[
 I_G(s)/s\ \longrightarrow\ J_G
 =\left|i\sum_j z_{Gj}Z_j\right|^2
\]
almost surely and in every fixed moment. Dominated convergence applies because
$I_G(s)/s\le n\|PZ\|^2$, which has every finite moment. The real Gaussian
field has covariance eigenvalues
\[
 \lambda_\pm=\frac{n\pm|M_2|}{2}.
\]
They are the same for A and B. Thus $J_A$ and $J_B$ have the same law as
$\lambda_+ Z_1^2+\lambda_- Z_2^2$ with independent standard normals.
Write $m_k=\mathbb E J_G^k>0$. Each exact intensity moment is a finite sum
of exponentials in $s$, hence is analytic at zero. Therefore, for every fixed
integer $k\ge1$,
\begin{equation}\label{eq:all-leading-moments}
 \mathbb E I_G(s)^k=m_k s^k+O(s^{k+1}),\qquad
 D_k(s)=O(s^{k+1}).
\end{equation}
The exact fourth-moment identity specifies the sharper coefficient
$D_4(s)=c_5s^5+O(s^6)$. Put $R=|M_2|^2=129187/400$. The generating function
\[
 \mathbb E e^{tJ_G}=
 \bigl[1-2nt+(n^2-R)t^2\bigr]^{-1/2}
\]
gives
\begin{equation}\label{eq:m4}
 m_4=24n^4+72n^2R+9R^2>0.
\end{equation}
These facts also justify $D_5=O(s^6)$ and $D_6=O(s^7)$ without assuming
an uncomputed sixth-order symbolic identity.

\subsection{Gaussian readout}

For fixed $\sigma>0$, the squared Hellinger distance between the laws of
$I_A+N(0,\sigma^2)$ and $I_B+N(0,\sigma^2)$ satisfies
\begin{equation}\label{eq:hellinger-gaussian}
 \boxed{\displaystyle
 h^2=\frac{c_5^2}{96\sigma^8}s^{10}+o(s^{10}).}
\end{equation}

\paragraph{Proof.}
Subtract the common mean $\mu(s)=\mathbb E I_G(s)=n(1-e^{-s})=O(s)$ from
both observations; this translation preserves Hellinger distance.
Set $V_G=I_G-\mu$. Then $\mathbb E V_G=0$, $|V_G|\le n^2$, and
$\mathbb E|V_G|^k=O(s^k)$ for every fixed $k$. Equality of the first three
raw moments gives equality of the first three central moments. The fourth
central-moment difference is $D_4$, and the fifth is
$D_5-5\mu D_4=O(s^6)$.

Use the Hilbert norm $\|f\|_{g}^{2}=\int |f(x)|^2/g_\sigma(x)\,dx$.
The translated Gaussian kernels obey the exact identity
\[
 \left\langle g_\sigma(\cdot-u),g_\sigma(\cdot-v)\right\rangle_g
 =e^{uv/\sigma^2}.
\]
It follows either by completing the square or direct Gaussian integration.
All fixed derivatives of these kernels have finite Hilbert norm, uniformly
for $|u|\le n^2$; the derivative norms are an exponential times a polynomial
in $u$. Taylor's theorem in this Hilbert space, through degree five, gives
\[
 \frac{p_{A,s}-p_{B,s}}{s^5}
 \longrightarrow \frac{c_5}{24}g_\sigma^{(4)}
 \quad\hbox{in }L^2(g_\sigma^{-1}dx).
\]
The sixth-order remainder is $O(s^6)$ in that norm by the uniform derivative
bound and $\mathbb E|V_A|^6+\mathbb E|V_B|^6=O(s^6)$.

Jensen's inequality supplies the needed denominator control globally:
\[
 \frac{p_{G,s}(x)}{g_\sigma(x)}
 =\mathbb E\exp\left(\frac{xV_G}{\sigma^2}
                 -\frac{V_G^2}{2\sigma^2}\right)
 \ge \exp\left(-\frac{\mathbb E V_G^2}{2\sigma^2}\right).
\]
This lower bound tends to one, and $p_{G,s}(x)/g_\sigma(x)\to1$ for each
fixed $x$. Hence
\[
 \frac{h^2}{s^{10}}
 =\int\frac{[(p_{A,s}-p_{B,s})/s^5]^2}
                 {(\sqrt{p_{A,s}}+\sqrt{p_{B,s}})^2}
 \longrightarrow
 \frac{c_5^2}{4\cdot24^2}
 \int\frac{(g_\sigma^{(4)})^2}{g_\sigma}.
\]
For completeness, the limit follows from Hilbert-norm convergence of the
numerator, which gives $L^1$ convergence of its weighted square; the extra
factor $g_\sigma/(\sqrt{p_A}+\sqrt{p_B})^2$ is uniformly bounded and tends
pointwise to $1/4$. Differentiating the kernel identity four times in each
variable at zero gives
$\int(g_\sigma^{(4)})^2/g_\sigma=4!/\sigma^8$.
Substitution proves \eqref{eq:hellinger-gaussian}.

\subsection{Photon counting with a fixed positive background}

For fixed $b>0$ and $\eta>0$ in $N\mid I\sim\operatorname{Poisson}(b+\eta I)$,
\begin{equation}\label{eq:hellinger-positive-background}
 \boxed{\displaystyle
 h^2=\frac{c_5^2\eta^8}{96b^4}s^{10}+o(s^{10}).}
\end{equation}

\paragraph{Proof.}
Let $p_\lambda(k)=e^{-\lambda}\lambda^k/k!$ and use the Hilbert norm
$\|f\|_b^2=\sum_k |f(k)|^2/p_b(k)$. Direct summation gives
\[
 \langle p_{b+u},p_{b+v}\rangle_b=e^{uv/b}.
\]
This identity ensures uniform Hilbert bounds for all fixed kernel derivatives
on the compact interval $0\le u\le\eta n^2$. Taylor expansion in intensity
through degree five, the exact lower-moment cancellations, and the sixth
moment bound give
\[
 (p_{A,s}-p_{B,s})/s^5
 \longrightarrow \frac{c_5\eta^4}{24}(S-\mathrm{Id})^4p_b
 \quad\hbox{in }\ell^2(p_b^{-1}).
\]
For every count $k$ and every admissible intensity,
\[
 \frac{p_{b+\eta I}(k)}{p_b(k)}
 =e^{-\eta I}(1+\eta I/b)^k\ge e^{-\eta n^2}.
\]
The mixture ratios tend pointwise to one and have the same positive lower
bound. The argument used for Gaussian readout therefore applies to the
Hellinger denominator. Differentiating the kernel identity gives
\[
 \sum_k\frac{[(S-\mathrm{Id})^4p_b(k)]^2}{p_b(k)}=\frac{4!}{b^4},
\]
which proves the stated coefficient. The assumption $b>0$ is essential;
this formula is not uniform as $b$ tends to zero.

\subsection{Zero-background counting and the four-count event}

At $b=0$ and fixed $\eta>0$ the scaling changes:
\begin{equation}\label{eq:hellinger-zero-background}
 \boxed{\displaystyle
 h^2=\frac{\eta^4c_5^2}{96m_4}s^6+O(s^7).}
\end{equation}

\paragraph{Proof.}
For each fixed count $k$,
\[
 p_{G,s}(k)=\frac{\eta^k}{k!}\mathbb E[I_G^k e^{-\eta I_G}]
 =\frac{\eta^km_k}{k!}s^k+O(s^{k+1}).
\]
For $k=0$, use $m_0=1$. Expanding the remaining exponential and using the
exact equality of moments of orders zero through three shows
$p_{A,s}(k)-p_{B,s}(k)=O(s^5)$ for $k=0,1,2,3$. Their Hellinger contributions
are respectively $O(s^{10}),O(s^9),O(s^8),O(s^7)$.
At $k=4$, the leading probability and the difference are
\[
 p_{G,s}(4)=\frac{\eta^4m_4}{24}s^4+O(s^5),\qquad
 p_{A,s}(4)-p_{B,s}(4)=\frac{\eta^4c_5}{24}s^5+O(s^6).
\]
The second formula uses $D_5=O(s^6)$ and the sixth-moment remainder.
Consequently
\[
 (\sqrt{p_{A,s}(4)}-\sqrt{p_{B,s}(4)})^2
 =\frac{\eta^4c_5^2}{96m_4}s^6+O(s^7).
\]
For $k=5,6$, equation \eqref{eq:all-leading-moments} yields probability
differences $O(s^6),O(s^7)$, so their contributions are $O(s^7),O(s^8)$.
Finally, the entire remaining tail is bounded without any nonuniform series:
\[
 \sum_{k\ge7}(\sqrt{p_{A,s}(k)}-\sqrt{p_{B,s}(k)})^2
 \le\Pr_A(N\ge7)+\Pr_B(N\ge7)
 \le\frac{\eta^7}{7!}(\mathbb E I_A^7+\mathbb E I_B^7)=O(s^7).
\]
Here $\mathbf1_{\{N\ge7\}}\le (N)_7/7!$ supplies the factorial-moment
bound. Summing all contributions proves \eqref{eq:hellinger-zero-background}.

Keeping only the binary indicator $\mathbf1_{\{N=4\}}$ gives the same
leading Hellinger coefficient: its complementary event contributes only
$O(s^{10})$. Thus a frequency estimate of exactly-four-count observations
achieves the optimal sample-order scaling in this limit. This concerns
repeated independent observations, not a claim that four-count events
dominate one-observation total variation at every noise level.

The earlier parity rule illustrates this distinction. Its even-count
probability is $1-\eta ns+O(s^2)$, and its A--B difference is
$(c_5\eta^4/3)s^5+O(s^6)$. If only parity is retained, the two Bernoulli
laws have
\[
 h^2_{\rm parity}=\frac{c_5^2\eta^7}{36n}s^9+O(s^{10}).
\]
Parity is asymptotically optimal for one observation in total variation,
but discarding count magnitude loses information for repeated measurements.

\subsection{From affinity to necessary and sufficient sample orders}

Let $m$ observations be independent conditional on the unknown structure,
with equal prior probabilities for A and B. The optimal success probability
is $(1+d_{\rm TV}(P_s^{\otimes m},Q_s^{\otimes m}))/2$. Direct Cauchy--Schwarz
and $\min(p,q)\le\sqrt{pq}$ give
\[
 1-\rho(s)^m\le d_{\rm TV}(P_s^{\otimes m},Q_s^{\otimes m})
 \le\sqrt{1-\rho(s)^{2m}}.
\]
These bounds and $\rho=1-h^2/2$ prove that, for any fixed target success
$a\in(1/2,1)$, the minimum sample count has order $1/h^2(s)$. More explicitly,
put $r=2a-1$. A necessary sample count is
\[
 m\ge\frac{\log(1-r^2)}{2\log\rho(s)},
\]
and taking the ceiling of $\log(1-r)/\log\rho(s)$ is sufficient for the
optimal likelihood-ratio test. Therefore the sharp orders are
\[
 \boxed{\begin{aligned}
  m_\mathrm{Gaussian}(s)&=\Theta(s^{-10}),&&\sigma>0\text{ fixed},\\
  m_\mathrm{Poisson}(s)&=\Theta(s^{-10}),&&b>0,\ \eta>0\text{ fixed},\\
  m_\mathrm{Poisson}(s)&=\Theta(s^{-6}),&&b=0,\ \eta>0\text{ fixed}.
 \end{aligned}}
\]
With only parity retained in the zero-background case, the sharp order
becomes $\Theta(s^{-9})$. The powers are proved; a universal finite-noise
sample number, an exact optimal prefactor for the minimum sample count,
and an experimental acquisition time are not. Calibration, independent
snapshots, fixed exposure/gain, and the assumed structural model remain
part of the theorem's conditions. Joint limits in $b$, $\eta$, $\sigma$,
or the observation frequency require a separate analysis.

